\documentclass[journal]{IEEEtran}
\usepackage{graphicx}
\usepackage{amsmath}
\usepackage{booktabs}
\usepackage{url}
\usepackage[UTF8,fontset=windows]{{ctex}}


\title{面向自动驾驶的车道感知：以人工锚定、引擎认证的真值为基础——显式计数契约与预注册单因子评测}
\author{袁哲宇 (Zheyu Yuan)}

\begin{document}
\maketitle
\begin{abstract}
自动驾驶中的车道感知被两个常被混为一谈的问题卡住：像素级人工标注的成本，以及「什么算一条车道线」缺乏公认定义。本文给出一个同时针对这两点的闭环研究计划。每个标签都带有一条显式的真值来源凭证------人工修订、智能体复核或引擎认证------因此任何结论都不必猜测自己立足在哪种真值上；所报告的轮次没有新增人工标注。评价本身被显式化并版本化：一份\emph{计数契约}为每个比率规定明确的分子与分母，一份\emph{协议}固定候选车道线的几何范围，并用哈希把封存的留出集绑定到它所依据的定义上。在两个参考集、一个按道路切分的留出集以及闭环驾驶中评测同一条分割臂（六个种子），得到三点发现。(i) 在权重固定的前提下，改变后处理开关与评价范围都会改变读数，甚至反转两条臂的排序；约 30\% 的标注线像素不满足采纳定义所用的外观谓词。(ii) 一次性留出确认\textbf{没有}复现验收数字（标签范围召回 0.460、精度 0.763）------这正是该独立检验的头条结果。(iii) 同一处定义变更伴随四次运行未观察到压中心线、车道来源可用率相对下降 66\%，而驾驶验收门仍为 0/4；残余阻塞是一处经测量确认的静止车体扫掠死锁。全部改动均预注册，并以事前固定的规则裁决；负结果与正结果以同等详细程度报告。
\end{abstract}

\section*{亮点}
- 真值来源凭证化并强制执行：人工修订、智能体复核、引擎认证三级。
- 计数契约为每一个报告的比率给出显式分子与分母。
- 协议哈希把封存的留出集绑定到它所依据的定义上。
- 「车道线」的两套自然定义在同一批数据上失败在不同的门上。
- 一次性留出确认未能复现开发集上的数字。
- 五个预注册驾驶层因子全部被否决；阻塞点是一处可测量的死锁。
\section{引言}
车道保持是自动驾驶栈最底层的能力，却仍然出奇地难以评测。文献中反复出现两个障碍。其一是标注成本：像素级车道标注昂贵，因此公开数据集相对于道路外观的多样性而言偏小，每进入一个新的运行域就要重新做一轮标注。其二是定义问题：论文报告的「车道线检测」精度，所依据的真值可能包含也可能排除褪色、被遮挡、虚线、黄色或处于阴影中的标线，也可能统计或不统计本车道之外的标线。于是两套系统可以在精度上相差几十个点，而它们唯一的分歧只是所检测对象的定义。
本文报告的是围绕这两个障碍建立起来的研究计划。它不是一个新的网络结构，而是试图把\emph{数据}与\emph{评价}做到足够严谨，使结构上的进展能够被度量。该计划包含三个组成部分。
\textbf{带凭证的真值来源。} 本计划中每一个标签都携带来源凭证，且凭证是被\emph{强制执行}而非仅仅声明：命令行的一句断言无法把一个标签「提升」到更高一级（见第 III-B 节）。凭证阶梯有三级------人工逐帧修订、逐帧智能体复核、引擎导出真值------而引擎一级只有通过四道显式认证门才被采纳：凭证门（标签来源必须声明且可核验）、漆面外观门（被声明为线的像素必须在图像中看起来像漆）、道路类型门（帧必须带有使「是否在铺装上」可判定的道路类型列）、负证据门（被声明为不含线的帧必须给出自身的正面证据，而不只是「没有标签」）。本文报告的各轮次不新增人工标注：实验变量是在人工修订基座之上、引擎认证增量的组成方式，第 VII 节说明基座规模偏小让我们付出了什么代价。
\textbf{显式计数契约。} 每一个报告的比率都由一份计数契约定义，契约在一条显式的集合链上指明分子与分母：含候选的帧、抽取出的候选实例、标注中存在的参考实例、匹配实例、横向位置一致的实例、左右角色一致的实例。因此覆盖率、身份率与角色一致率是分母明确的实例级比率；精度与召回仍是像素级，并以两种\emph{范围}（标签像素、类漆像素）报告，因为二者回答的是不同的问题。
\textbf{带哈希门控的版本化协议定义。} 协议固定候选车道线的几何范围，本身就是一等、可版本化的产物。协议变更会改变评价哈希，而被冻结的留出集与一次性确认都记录该哈希，因此在一套定义下产生的确认不可能被悄悄归到另一套定义名下。该机制在工作中捕获了一个真实缺陷：在修复之前，哈希并不跟随协议，于是按旧定义封存的集合将能授权一次新定义下的确认。
\textbf{方法纪律。} 本文的每一处改动都在运行之前预注册、并写明读数规则，且每份预注册都标明它所冻结的协议版本与仓库提交；规则包括：安全非劣性按各次运行的最大值读取、明确拒绝把未测量的量当作通过、禁止放宽已冻结的门。改动按单因子逐个施加，并以交替的对照/因子运行裁决。当预注册规则给出的裁决我们认为过于粗糙时，我们照原样报告裁决，并记录规则本身的缺陷，而不是重新解释结果；其中一处缺陷------安全指标是帧\emph{计数}、因而随车辆行驶距离缩放------在第 V-F 节分析，并在补充材料中量化。
\textbf{贡献。} （1）带凭证的真值来源阶梯（人工修订／智能体复核／引擎认证）及其强制执行，以及针对机器生成增量的认证门流水线；每个数据池的组成都被显式说明，且所报告的轮次不新增人工标注。（2）计数契约，以及「线」的两套自然读法在同一批数据上失败在不同门上的实证------它把约 30\% 的标注线像素移出采纳范围（采纳范围保留开发池 69\%、有限类别池 89\% 的标注线像素）。（3）协议／版本／哈希机制，把留出集绑定到它所依据的定义上。（4）一次按道路切分的一次性确认，配有显式的 UNKNOWN 政策，其负结果（未复现开发集数字）被我们当作发现来报告。（5）对「为什么一套通过全部离线门的感知栈仍然过不了闭环验收门」的测量式定位，包括一个静止死锁机制，以及「在铺装上」门实际上是由神经障碍层而非铺装证据驱动的这一观察。
\section{相关工作}
\textbf{车道线检测方法。} 现代车道线检测由深度分割与参数化／锚点式曲线回归主导。带空洞卷积与判别损失的编码器--解码器分割（LaneNet \cite{neven2018lanenet}）与逐行消息传递（SCNN \cite{pan2018scnn}）确立了本文系统所沿袭的分割路线；基于锚点头的跨层细化（CLRNet \cite{zheng2022clrnet}）提升了曲线精度与速度。同一族的中文工作探索了 CLRNet 的轻量化变体 \cite{chen2026clrnetlight}、B 样条端到端回归 \cite{huang2026bspline}、图／Transformer 混合模型 \cite{jia2026gnn}、带交叉注意力的 Transformer \cite{li2026xca}、图像序列上的并行检测 \cite{zhu2024parallel}、多尺度特征聚合 \cite{hu2026multiscale}、DeeplabV3+ 变体 \cite{shuai2026deeplab}，其原始架构见 \cite{chen2018deeplab}、面向复杂场景的对称先验与关键点方法 \cite{lu2026symmetry}，以及基于卷积架构的 3D 车道线估计 \cite{chen2025threed}。这些工作与本文研究对象相同，但优化目标不同：它们报告的精度所依据的数据集，其标注约定彼此并不一致------而这正是本文要显式化的那种歧义。
\textbf{数据集与真值。} TuSimple \cite{tusimple2017}、CULane \cite{pan2018scnn} 与 Cityscapes \cite{cordts2016cityscapes} 仍是标准基准；二者的标注约定（逐行车道存在性、远场标线是否计入、遮挡的处理方式）不同，且都没有报告标注的\emph{外观}属性。真值生成本身已被作为独立问题研究：Borkar 等 \cite{borkar2012groundtruth} 为夜间车道检测器构建了基于时间切片可视化的自动真值流程，McCall 与 Trivedi \cite{mccall2006video} 则明确从评测方法学的角度综述并评估了基于视频的车道估计。本文延续这一路线并补充两点：真值来自引擎状态而非图像；以及一份使每个报告比率的分母都可被审查的计数契约。
\textbf{自动标注与仿真。} 自训练与自动标注被广泛用于降低标注成本，但其评价继承了上述定义问题。由仿真器导出的标签在合成基准中很常见；不那么常见的是对每个生成标签针对图像证据做认证的纪律------本文采用它，因为一个无法从图像中找回的标签不是可学习的目标。
\textbf{评价纪律。} 预注册在实证科学中是标准做法，在感知研究中却很罕见。我们在此的贡献不是概念，而是一个完整记录的实例：一次预注册、一次交替的单因子 A/B、一条固定的读数规则，以及每个因子一份裁决------其中五个因子的裁决为负。
\section{系统与方法}
\subsection{总体架构}
本栈是一个以感知为主导的驾驶系统。一个相机环（八个安装位）与一份激光雷达距离图像送入语义分割头，其线输出被转换为世界坐标下的标线；在协议定义下从这些标线中选出一条车道参考；规划器生成路径；安全监视器按硬检查（位姿、边界、占用、障碍风险）对每个节拍进行仲裁，并发布速度与等级；控制器消费仲裁后的路径。本栈的铁律是：横向参考只能来自感知------地图只回答「去哪里」，从不回答「车道在哪里」。补充材料给出候选范围流水线，并标注了第 V-E 节讨论的实测撤销率。
\subsection{真值来源、凭证与认证门}
这里必须区分两个维度。\textbf{来源}说的是标签由谁／什么产生（人工修订、智能体复核、引擎）；\textbf{认证}说的是该标签是否通过下面四道门。人工修订的帧可以不过外观门，引擎帧也可以四门全过，因此准入由门的结论决定、而不是由级别决定。在这个前提下，本计划区分三级真值，并为每个标签记录其级别：\textbf{人工修订}（由人逐帧标注或修正车道标签）、\textbf{智能体复核}（对引擎输出的逐帧智能体审查）、\textbf{引擎认证}真值（引擎自身的车道几何投影到相机坐标系）。有两个池是人工修订的：训练基座（5 个目录、75 帧）与开发／评价包（9 个包、117 帧），这也是第 V-A 节的验收结果是对照人工标注测量的原因。消融实验所变动的机器生成增量是引擎认证的；它只有在满足以下条件时才被采纳：(i) 凭证存在且可解析，且标签来源已声明（来源无法命名的帧被排除，而不是猜测）；(ii) 被声明为线的像素满足漆面外观谓词------白漆要求高亮度、低色度与红蓝平衡；黄漆要求特定的色调比区间；该谓词是单一来源模块，因此阈值变更是一次版本化变更，而不是局部修改；(iii) 帧带有道路类型列，缺少它则「在铺装上」不可判定；(iv) 被声明为不含线的帧携带自身的正面证据（已认证负例），因此「没有标签」永远不会被读成「没有线」。任何一道门不通过的帧都从训练与评价中排除，并计入审计，而不是被悄悄丢弃。
门的记账本身就是一项结果：补充材料给出某一采集轮次的负例池漏斗（采纳、可疑、含漆、未决）------一个在磁盘上看起来很充裕的池，在被用作训练信号之前就收缩为一个已认证子集。
\subsection{计数契约}
契约定义在一组逐帧累加的整数计数之上，比率由这些整数一次算出，绝不对各目录的比率再取平均。$P$ 是标注在铺装上确实有线的帧数（逐帧判定：不因同组另一帧有线而纳入）；$C$ 是 $P$ 帧内进入评价的感知候选数；$C_{\mathrm{out}}$ 是落在 $P$ 之外帧上的候选，单独上报；$R \subseteq C$ 是自身所在侧有可用真值参考的候选数；$M \subseteq R$ 是满足冻结匹配条件的候选数；$L \subseteq M$ 是预测与参考都有可判左右角色的候选数；$A \subseteq L$ 是角色一致的候选数。报告的实例级比率为覆盖率 $|R|/|C|$、身份率 $|M|/|R|$ 与角色一致率 $|A|/|L|$；像素级精度与召回以两种范围（标签、漆）报告。图 1 给出该契约。
\begin{figure*}[!t]
\centering
\includegraphics[width=\textwidth]{figures_zh/fig30_counting_contract.png}
\caption{计数契约（示意）：每个比率都在冻结的计数定义上有显式分子与分母；覆盖率是候选参考可判率，不是检出率。}
\label{fig:contract}
\end{figure*}
有两条推论必须从定义读出，而不能想当然。第一，覆盖率是\textbf{候选参考可判率}：进入评价的候选中「自身所在侧有真值可比」的比例，而不是「真实标线被检出」的比例。一个有线但感知没有产生任何候选的帧不向 $|R|/|C|$ 贡献任何东西，它的漏检由像素级召回承担。第二，落在 $P$ 之外的候选记为 $C_{\mathrm{out}}$ 并单独上报，因此「没有参考可判」永远不会被算成覆盖率失败。据此，本文\textbf{不主张实例级召回}：上面三个条件比率不能替代它；要报告实例召回，必须给出真值实例总数与显式的一对一分配规则。每个比率仍有显式分母，且帧级计数 $P$ 与候选级计数分开报告，因此「含候选帧很少」的运行无法伪装成一次好运行。
\subsection{协议定义}
协议固定候选车道线的三件事：其\textbf{横向范围}（离本车所在车道超过阈值的候选不算「我车道的线」）、\textbf{同侧近重复标线是否合并}为一个实例、以及\textbf{外观门是否剔除}像素看起来不像漆的候选。冻结的基线协议三项全不启用（横向范围关、合并关、外观关）；采纳协议三项全开，横向阈值 5.5 m、合并半径 1.0 m，外观谓词见第 III-B 节。补充材料给出开关矩阵与实时哈希；哈希是当前启用协议的函数，因此评价产物不可能被跨定义改属。
横向范围并非任意取值：它由参考池的实测几何设定，见补充材料（261 042 个标线像素，中位偏移 −0.17 m，中位绝对偏移 1.22 m）。5.5 m 的范围保留本车道标线及其紧邻标线，同时排除主导了第 V-E 节所讨论错误配对的远场标线。
\subsection{训练组成的设计空间}
由于真值是生成而非标注而来，训练池的组成本身就是一等的实验变量。我们沿一个轴研究它：\textbf{负例剂量}，即相对于基础组成混入多少已认证的无标线帧，以及把标线置于特定几何关系（近／远配对、同侧配对）的若干组成变体。第 V-C 节报告由此得到的剂量--效应。
\subsection{留出最终集：封存、哈希门控、一次性消费}
独立确认集按三条约束组成：它必须与所有训练池\textbf{道路不相交}（身份取自生成元数据，而非目录名）；它必须通过第 III-B 节各门认证；其内容必须由逐帧哈希的摘要冻结。封存记录帧摘要、协议哈希与组成；访问随后以协议哈希为门控，且一次确认会\textbf{消费}该集合：确认成功后该集合不得再被读取，被拒绝的访问只记入台账而不消费。确认程序在封存帧上评测候选，并记录哪些量被测量；未被测量的量记为 UNKNOWN，不得被报告为已确认。完整流程与实时封存信息见补充材料。
\subsection{驾驶栈与安全仲裁}
在闭环中，同一套感知驱动规划器与安全监视器；监视器的仲裁阶梯按固定顺序评估规则；软规则累积速度上限，而最坏等级为 `minimal_risk` 的规则会停止该节拍。补充材料列出该阶梯。第 V-E 节涉及其中两条规则：\textbf{规划车体扫掠规则}，当规划路径的车体扫掠会以小于硬阈值（4 m）的距离穿越车道边界时停止该节拍；以及\textbf{铺装门}，当车道中心不落在观测到的可行驶表面上时撤销感知车道。
\section{实验设置}
\textbf{参考集。} 主验收使用两个不相交集合：由 9 个已认证场景组成的开发池（117 帧），以及为近场背景结构而挑选的有限类别池（每臂 24 帧）。二者都不与训练池共享帧。
有限类别池是被\emph{挑选}而非抽样的：候选场景按一个与模型无关的近场结构密度准则排序，该准则由人工标注计算得到（补充材料），因此该池集中在身份与角色最难的场景上，而不是平均场景上。第三个按道路切分的集合共 200 帧（其中 96 帧含线），被封存用于第 V-D 节的一次性确认。
不相交性用位姿审计，而不是用名字：每个候选采集都在 50 m 缓冲下与开发锚点比对，并留下一条已记录的裁决（见补充材料）。同一套机制曾否决封存集的前两个组成方案，这也是封存组成另带一份内容级审计的原因（补充材料）。
\textbf{实验臂。} 所有分割臂共享同一初始化、同一预算与同一配方；它们只在训练池组成（剂量臂）或单一后处理因子（协议组件）上不同。组成是显式的：5 个目录（75 帧）的人工修订基座，加上数量可变的引擎认证包（所报告基础配置中为 2 个配对线包共 16 帧、12 个已认证负例包共 96 帧），并以负例剂量作为变动轴。验收使用六个随机种子；组成臂按各自的种子数报告。
\textbf{指标与门。} 冻结的门为覆盖率 $\ge 0.80$、身份率 $\ge 0.60$、精度 $\ge 0.40$、召回 $\ge 0.70$、角色一致率 $\ge 0.70$；召回对采纳协议按漆范围读取，对基线按标签范围读取，且两种读法始终同时报告。两轮的裁决规则并不相同，而且这一差别本身就是记录的一部分：基线轮的存储判定要求\textbf{每个种子都通过}，采纳轮则显式记录 `verdict_mode = mean` 并同时给出逐种子数值。因此第 V-A 节既给出各自规则下的存储判定，也把同一批逐种子数值在两种规则下重新汇总，使任何一个门的结论都不必在未言明的规则下被解读。
\textbf{驾驶实验。} 闭环运行使用一个固定场景（一段城镇路段），交替设置对照臂与因子臂，每臂四次运行；硬目标即基准的固定检查（帧数、倒车、按参考线与按车体扫掠的中心线／边缘穿越、停滞、在路、碰撞、到达目标）。读数规则在每一轮之前固定：硬门要求 4/4，安全项按各次运行的最大值比较，可用率按中位数与范围比较，UNKNOWN 永不解除门。
\textbf{纪律。} 第 V 节的每个因子都带读数规则预注册；全程未放宽任何阈值；
计时是记录而非估计：补充材料是任何延迟声明之前所用的重复测量协议（安静机器前置条件、重复运行、取较小值------曾有一个泄漏的仿真器实例把 p95 抬高到 7.6 倍）。每份判定文件、清单与逐帧遥测都被保留，本文每张图都由这些文件经脚本重新生成。
\section{结果}
\subsection{两套定义下的验收}
表 I 在两个集合上报告两套定义，表 II 则把同一批检查点的召回门在后处理 $\times$ 范围两种设置下重新汇总。在基线定义下，存储判定在两个集合上都失败身份门（开发 0.433，有限类别 0.414），有限类别池还额外失败精度（0.402）与角色（0.674）；在采纳定义下，按均值规则五个门在两个集合上全部通过（覆盖率 0.956/0.861、身份率 0.665/0.716、精度 0.894/0.659、漆范围召回 0.812/0.885、角色 0.883/0.793）。图 2 把同样的数字画在阈值之上。
\begin{figure*}[!t]
\centering
\includegraphics[width=\textwidth]{figures_zh/fig1_gate_matrix.png}
\caption{两套协议定义下的门指标与冻结阈值：开发池 (a) 与有限类别池 (b)。}
\label{fig:gates}
\end{figure*}
这张表要带三条限定。(i) 两行并非在同一规则下裁决。若统一按均值规则重算，基线在有限类别池上的精度（0.4018）高于 0.40 阈值，尽管 6 个种子里有 4 个低于它；基线在开发池上的召回（0.736）高于 0.70 阈值，尽管 6 个种子里有 2 个低于它。若统一按「全部种子通过」规则重算，采纳定义的开发池标签范围召回 6/6 种子不通过，而其漆范围召回只有 1/6 不通过。(ii) 协议变更同时改动了三个后处理开关与评价范围，因此这张表不能把差异整体归因于定义；表 II 把范围与后处理在召回上分开，逐个开关的单因子结果见第 V-B 节。(iii) 候选级门是图景中较稳定的部分：基线定义的开发池与有限类别池各有 6/6 种子失败身份门，有限类别池另有 5/6 种子失败角色门；而在采纳定义下没有任何种子在这三个门上失败（补充材料）。
补充材料在两套定义下逐种子给出有限类别池，并标出未达标的候选级门；它与表 I 取自同一批判定文件，因此图与表不可能对「数字来自哪一次运行」产生分歧。
\begin{table*}[!t]
\centering
\caption{两套协议定义在两个参考集上的结果，按各轮自己的裁决规则存储（6 种子均值；加粗表示存储判定未通过的门；采纳协议按漆范围读取，基线按标签范围读取）。两行的规则并不相同------见正文与表 II。}
\begin{tabular}{@{}lllllll@{}}
\toprule
集合 & 定义 & 覆盖率 & 身份率 & 精度 & 召回 & 角色 \\
\midrule
开发（9 场景，117 帧） & 基线（标签范围） & 0.952 & \textbf{0.433} & 0.610 & 0.736 & 0.868 \\
开发 & 采纳（漆范围） & 0.956 & 0.665 & 0.894 & 0.812 & 0.883 \\
有限类别（24 帧） & 基线（标签范围） & 0.819 & \textbf{0.414} & \textbf{0.402} & 0.804 & \textbf{0.674} \\
有限类别 & 采纳（漆范围） & 0.861 & 0.716 & 0.659 & 0.885 & 0.793 \\
\bottomrule
\end{tabular}
\end{table*}
\begin{table}[!t]
\centering
\caption{召回门的重新汇总：每个池的同一批 6 个检查点，在后处理 $\times$ 范围两种设置下读取，并给出低于 0.70 阈值的种子数。漆范围按采纳定义读取该检查点，标签范围按基线定义读取。}
\begin{tabular}{@{}lllll@{}}
\toprule
池 & 后处理 & 范围 & 召回均值 & 低于 0.70 的种子 \\
\midrule
开发 & 基线开关（v7） & 标签 & 0.736 & 2/6 \\
开发 & 基线开关（v7） & 漆 & 0.813 & 1/6 \\
开发 & 采纳开关（v8） & 标签 & 0.559 & 6/6 \\
开发 & 采纳开关（v8） & 漆 & 0.812 & 1/6 \\
有限类别 & 基线开关（v7） & 标签 & 0.804 & 0/6 \\
有限类别 & 基线开关（v7） & 漆 & 0.891 & 0/6 \\
有限类别 & 采纳开关（v8） & 标签 & 0.775 & 0/6 \\
有限类别 & 采纳开关（v8） & 漆 & 0.885 & 0/6 \\
\bottomrule
\end{tabular}
\end{table}
\subsection{分歧是定义性的}
这里同时有两件事在动，必须分开。\textbf{范围}改变读数：表 II 显示在后处理固定的前提下，从标签范围换到漆范围，开发池召回由 0.736 变为 0.813（基线开关）、由 0.559 变为 0.812（采纳开关）------后一种情形是在权重毫无变化的情况下出现 25 点的摆动。\textbf{后处理}同样改变读数，而且方向不同：把范围固定在标签，采纳开关付出 18 点召回（0.736 $\rightarrow$ 0.559）。表 V 把每个开关单独打开，把这份代价完全归给外观门：横向范围与同侧合并都不动标签召回（两者都是 0.736），它们动的是身份率（分别 0.432 $\rightarrow$ 0.621 与 0.432 $\rightarrow$ 0.459）；外观门是唯一用标签召回换像素精度的开关。范围变更在像素上值多少是量出来的、不是假设的：漆范围保留开发池 69\%、有限类别池 89\% 的标注线像素；补充材料逐种子给出两种范围下的召回，并叠加标签线像素中的非漆比例（约 0.30）。这些被排除的像素正是不满足外观谓词的那一批；第 VI 节说明这一测量允许我们主张什么、不允许主张什么。
实例级门比这更敏感，而且敏感的方式直接关系到选模。补充材料在标签范围与表面范围两种身份定义下读取 20 个检查点（两条臂、每臂 10 个种子）：基础臂均值由 0.399 变为 0.867，6$\times$ 臂由 0.429 变为 0.771。一条臂的绝对值几乎翻倍、另一条反而下降，而且\textbf{两条臂的排序发生反转}（标签范围下基础臂低于 6$\times$ 臂 0.031，表面范围下高于它 0.095）。同一处反转记录在仓库关于表面范围的协议提案里，这也是该范围只上报、不升为主口径的原因：一个能改变选模结果的阈值是一项选模决策，而不是报告细节。因此身份阈值必须像召回门一样指明其范围。该实验使用自己的池与每臂 10 个种子，其数字与表 I 的 6 种子均值不可比，本文也不作这样的比较。
这一效应并不只属于交付臂。在一个用不同组成训练的第二个臂上（近／远配对包、六种子、同一冻结协议），从标签范围换到漆范围把召回由 0.459 提到 0.673（+0.214），交付臂为 +0.252；身份门读作 0.658 对 0.665。范围效应的方向在第二个模型上复现，而幅度依臂而变------这正是本文把范围与每个数字一起报告、而不是给出一个统一校准值的原因。
另有两项单因子结果约束了采纳定义的行为。补充材料（边界图）扫描掩码膨胀：膨胀在两个集合上都换来召回、付出精度，而有限类别池为此付出的代价大得多（膨胀一个像素时其精度由 0.65 降到 0.39，降幅 0.26；开发池由 0.89 降到 0.84，降幅 0.05；两行各按自己的膨胀档位读取），因此没有任何单一膨胀值能同时通过两个集合。图 3 显示外观门完全按设计起作用：精度显著上升（开发 0.650 $\rightarrow$ 0.884；有限类别 0.387 $\rightarrow$ 0.650），有限类别池的 IoU 改善（0.347 $\rightarrow$ 0.532），代价是召回（开发 0.766 $\rightarrow$ 0.595）。
\begin{figure*}[!t]
\centering
\includegraphics[width=\textwidth]{figures_zh/fig14_appearance_gate.png}
\caption{外观门对像素指标（线 IoU、精度、召回）的消融，两个集合，种子 42------它是唯一用标签召回换精度的开关（表 V）。}
\label{fig:appearance}
\end{figure*}
外观谓词是掩码门，不是真值判据：该模块自身记录了它在已认证的无标线帧上有 13k--68k px 的命中（亮铺装与亮碎石都能满足白色分支），这也是它被用来剔除候选、而不是用来宣告「哪里有标线」的原因。在这个前提下，实例级探针解释了外观门为何可以足够激进：在已认证的有限类别场景中，参考实例被标注覆盖的程度基本为 1.0，且与横向偏移无关，因此存活下来的候选仍然是对照完整真值被测量的。另一处后处理杠杆是并行合并半径，它移除同侧近重复候选（半径从 6.0 m 收紧到 3.0 m 时，开发池保留数 477 $\rightarrow$ 166），代价是身份率------这也是采纳半径取 1.0 m 而非扫描激进端的原因。两个掩码阈值同样被扫描过；在该池上负例一侧不可测（没有合格负例帧），因此该扫描报告在补充材料中，而不是画成零。
\begin{table*}[!t]
\centering
\caption{冻结开发池上的后处理单因子，逐个开关单独开启（每格 6 种子；C0 为基线定义，C4 为采纳定义）。横向范围是身份率杠杆（0.432 $\rightarrow$ 0.621，六个种子全部为正）；外观门是精度杠杆，以标签召回为代价（精度 0.610 $\rightarrow$ 0.894、召回 0.736 $\rightarrow$ 0.559），且不提升身份率；合并是小的身份率效应（+0.027）。}
\begin{tabular}{@{}llllll@{}}
\toprule
配置 & 标签精度 & 标签召回 & 漆召回 & 真值像素保留 & 身份率 \\
\midrule
v7：三项全关 & 0.610 & 0.736 & 0.813 & 0.693 & 0.432 \\
+ 横向范围 5.5 m & 0.610 & 0.736 & 0.812 & 0.693 & 0.621 \\
+ 同侧合并 & 0.610 & 0.736 & 0.813 & 0.693 & 0.459 \\
+ 外观门 & 0.894 & 0.559 & 0.812 & 0.693 & 0.409 \\
v8：三项全开（采纳） & 0.894 & 0.559 & 0.812 & 0.693 & 0.665 \\
\bottomrule
\end{tabular}
\end{table*}
\subsection{训练组成的剂量--效应}
图 4 给出剂量响应当前的样子，而旧曲线没有通过它自己的检验。旧曲线（基础 1 种子、4$\times$ 6、6$\times$ 1、8.5$\times$ 6：0.707/0.711/\textbf{0.763}/0.701）是在另一套配置下测的，其身份率在冻结协议下无法复现------同一个 6$\times$ 检查点在那里读作 0.763、在这里读作 0.667------因此它与验收数字不可比，已被替换。新图报告冻结协议下的\textbf{配对种子重评}（凡检查点仍在的剂量都用种子 42--47；见表 VI）：0$\times$ 0.658（0.631--0.690）、4$\times$ 0.661（0.648--0.674）、6$\times$ 0.665（0.645--0.676）。相对 0$\times$ 的逐种子配对差值为 +0.003（4$\times$）与 +0.007（6$\times$），且符号不一致（两者在种子 45 上都是负的），因此\textbf{六种子下三档不可区分}，「6$\times$ 是剂量响应峰值」这一主张不成立。8.5$\times$ 臂的检查点已在存储清理中删除，其\emph{已记录}六种子均值为 0.701------高于 6$\times$，与旧曲线的排序相反。交付组成仍是 6$\times$，因为它在一次性确认之前就被选定，改动它需要重新封存一个集合；论文对组成的主张现在被限定为「未证明存在剂量效应」。
\begin{figure*}[!t]
\centering
\includegraphics[width=\textwidth]{figures_zh/fig3_dose_response.png}
\caption{冻结协议下配对种子（42--47）的剂量响应：重评的三档不可区分，而已记录的 8.5$\times$ 值更高（表 VI）。}
\label{fig:dose}
\end{figure*}
本计划的训练历史（44 个家族、321 次已记录运行，见补充材料）在 120 轮预算的 20--40 轮内收敛，而线通道被掩蔽的家族按构造即处于零线 IoU（图中已标注），因此它们的曲线不是负结果。补充材料转而扫描横向范围：从 6.0 m 收紧到 3.0 m 提升身份率（开发 0.582 $\rightarrow$ 0.713；有限类别 0.581 $\rightarrow$ 0.758）并减少候选（477 $\rightarrow$ 166 与 166 $\rightarrow$ 127），而角色一致率持平------横向范围是身份杠杆，不是角色杠杆。另有两个单因子被否决：在最后五个轮次上做随机权重平均 \cite{izmailov2018swa} 没有救回任何种子；一个 Tversky 损失权重臂把负例一侧的假阳性帧率放大到 4.5 倍（占合格帧 11.5\% $\rightarrow$ 51.9\%，即 +40.4 个百分点；假阳性像素 329 $\rightarrow$ 1730；最大连通域 320 $\rightarrow$ 808 px），这也是交付臂保留平衡损失的原因。
在所有已记录的像素臂上，负例一侧的读数呈聚集而非发散（见补充材料），这也是负例剂量必须被当作变动轴、而不是去调损失权重旋钮的原因。把剂量--效应扩展到五个指标（见补充材料）显示：身份率与角色随剂量变动，而覆盖率与路外候选比例几乎不动。采纳定义的可用率代价在配对阶段本身也可见：每 117 帧运行损失 3--4 个配对帧。
\begin{table}[!t]
\centering
\caption{冻结协议下的等种子剂量对照（各档配对使用种子 42--47）。重评的三档不可区分；已记录的 8.5$\times$ 值高于 6$\times$，因此此前的「6$\times$ 是峰值」不成立。}
\begin{tabular}{@{}lllll@{}}
\toprule
剂量 & 身份率均值 & 身份率范围 & 相对 0$\times$ 的配对差值 & 角色 \\
\midrule
0$\times$（基础） & 0.658 & 0.631--0.690 & --- & 0.874 \\
4$\times$ & 0.661 & 0.648--0.674 & +0.0030 & 0.864 \\
6$\times$（交付） & 0.665 & 0.645--0.676 & +0.0069 & 0.883 \\
8.5$\times$（已记录；检查点已删除） & 0.701 & 0.655--0.744 & --- & --- \\
\bottomrule
\end{tabular}
\end{table}
\subsection{一次性留出确认}
封存集被读取过一次，读取者是交付候选。总体标签范围召回为 \textbf{0.460}、精度 \textbf{0.763}（IoU 0.403；道路 IoU 0.633）；分组来看，七个分组中有三组召回落在 0.17--0.20，但精度并不一致：其中两组精度为 0.864 与 0.922，第三组（a27）精度只有 \textbf{0.169}，其预测的线像素有 61\% 落在路外------这一组是候选基本失败的分组，而不是一个干净的高精度分组。有一个分组不含线真值（召回无定义，已按无定义记录）；其余分组落在 0.43--0.63。负例一侧，105 个合格帧中有 22 帧带假阳性线像素（20.9\%），但这些像素仅占合格像素的 0.084\%，最大连通域为 1164 px，且控制区域内没有出现任何假候选。图 5 给出该确认。这一结果有两点值得强调。第一，它\textbf{没有}复现第 V-A 节的验收数字：独立集更难，其分组组成也不同。组成有两种计法，两种都记录在案：按包计，200 帧中 104 帧来自已认证的无标线包、96 帧来自含线包；按逐帧资格计，95 帧含线真值、105 帧为合格负例------两者的差别来自「帧自身资格」与「包标签」不一致的那些帧。交付候选的主张以上文所报数字为界，而不是以验收表为界。第二，确认记录中明确写明实例级探针（身份率、角色、参考覆盖）\textbf{未}在该集上运行；它们是 UNKNOWN，不得被报告为已确认。
\begin{figure*}[!t]
\centering
\includegraphics[width=\textwidth]{figures_zh/fig4_final_confirm.png}
\caption{封存的道路不相交集上的一次性最终确认：按组指标 (a) 与负例一侧诊断 (b)。}
\label{fig:confirm}
\end{figure*}
补充材料说明确认所运行的组成（七个分组共 200 帧；摘要与协议哈希印在标题中，因此该图不可能被改属到另一次封存）。补充材料是准入它的那次审计：每个来源包都与已使用的 407 个帧哈希比对，出现重叠的包在封存前被剔除。第三点与第 V-A 节的差距有关：封存集的真值是\textbf{引擎认证}的，而验收集是\textbf{人工修订}的------两个数字不仅对照不同数据，也对照不同真值级别。
\subsection{闭环}
\textbf{闭环中的定义变更。} 在基线定义下，闭环运行的中心线穿越为 8、8、8、13 帧；在采纳定义下，四次运行均未观察到压中心线。这是关于四次运行的陈述，不是一个率：同一批运行行进极少（停滞比例 0.93--1.00，见下），在接近零的暴露量下，零事件不能被外推。采纳定义为这一点付出可用率代价：传感器车道来源率从 0.851（范围 0.231--0.975）的中位数降至 0.292（0.097--0.524），相对下降 66\%；配对车道率从 0.890（0.809--0.986）降至 0.528（0.130--0.995）。图 6 给出该权衡。两种效应都与离线发现一致：基线保留远场与非漆候选，而这些候选的配对会把车辆拉过车道中心。
\begin{figure*}[!t]
\centering
\includegraphics[width=\textwidth]{figures_zh/fig8_closed_loop_tradeoff.png}
\caption{闭环权衡：中心线穿越与传感器车道来源率，每臂四次运行。}
\label{fig:closedloop}
\end{figure*}
\textbf{验收失败及其机制。} 驾驶验收门（到达目标＋感知主导横向模式＋严格仲裁，四次运行）为 \textbf{0/4}：停滞比例 0.93--1.00，120 s 内行进 0.9--14.5 m，距目标剩余 76.5--88.8 m；已测量的安全项保持（0 碰撞、0 路外帧、0 倒车）。补充材料给出最近各次运行的逐运行硬目标矩阵。随后我们以测量而非猜测定位了阻塞点。在 12 次城镇运行中，原因为「规划车体穿越车道边界」的帧共 512 个；在这些帧中，96\% 的情况下车辆静止，当前车体在 512/512 帧中都位于车道内，规划穿越距离的中位数为 2.50 m（509/512 落在 2--4 m）------即恰好位于 4 m 硬停阈值之内。车辆停在距车道中心 0.86 m 处，航向偏离车道走向约 11°，因此规划器实际生成的路径，其膨胀车体扫掠会在约 2.5 m 内穿越边界。记录能支持的是：\emph{已生成}的路径被反复拒绝，且在该集合的每一帧里都是如此；它不能证明不存在动力学上可行的前向路径，因为规划器的候选集并不是全部可行路径的集合。在冻结位姿上做受控扫描列在未完成工作中。图 7 给出该解剖。
逐运行轨迹（速度与路径到车道偏差）、各臂分布与车道放置分布见补充材料；同批运行中测量的安全裕度（最近障碍、最小碰撞时间、路径占用）是「验收失败属于可用率失败而非安全失败」这一结论的依据。
\textbf{四个因子已运行，一个提议在实现前被否定。} 针对该阻塞点我们预注册了五个驾驶层因子；其中四个被实现并运行，第五个（铺装重定中）在写任何代码之前就被既有测量否定（表 III）。一个让已验证轨迹在其提议过期后仍可复用的有界「保持」窗口提升了运动量（复用 7 $\rightarrow$ 42 帧；行进 3.15 $\rightarrow$ 8.5 m），但恶化了最坏情况的车体穿越计数（29 $\rightarrow$ 30）。把规划路径拉到感知车道中心在 39 帧上生效，并显著改善了最坏情况（中心线穿越 3 $\rightarrow$ 0；车体穿越 84 $\rightarrow$ 39），但恶化了一项右侧指标（0 $\rightarrow$ 2）。从铺装门中去掉障碍项（见下文）使行进距离翻倍（6.85 $\rightarrow$ 14.1 m）、距目标距离减半（83.5 $\rightarrow$ 76.6 m），但在车辆行驶两倍距离的同时恶化了车体穿越（23 $\rightarrow$ 66）。换用交付感知臂在路径可用率上没有产生分离（0.179 $\rightarrow$ 0.200，区间重叠），且最坏情况更差（36 $\rightarrow$ 49）。四个已实现的因子全部遵循预注册的最大值规则，并作为负结果报告；没有任何因子被采纳，第五个按构造仍未测试。
仲裁直方图是停滞比例背后的证据：在 2443 个已结算帧（运行开始至少 8 s 的帧，汇总自图中点名的十二次城镇运行）中，「无可用路径」是单项最大原因，为 \textbf{1383 帧}；其次是规划车体扫掠穿越 \textbf{512 帧}，二者合计构成静止时间的主要部分。
\begin{table*}[!t]
\centering
\caption{针对死锁的预注册驾驶层因子：机制是否生效、运动读数与预注册裁决。四个已实现并运行；第五个在实现前被既有测量否定。}
\begin{tabular}{@{}llllll@{}}
\toprule
因子 & 机制是否生效 & 行进（中位） & 停滞（中位） & 最坏车体穿越帧（对照 $\rightarrow$ 因子） & 裁决 \\
\midrule
有界保持窗口 & 是（复用 7 $\rightarrow$ 42 帧） & 3.15 $\rightarrow$ 8.5 m & 0.986 $\rightarrow$ 0.958 & 29 $\rightarrow$ 30 & 否决 \\
路径重定中 & 是（39 帧） & 5.1 $\rightarrow$ 5.75 m & 0.984 $\rightarrow$ 0.980 & 84 $\rightarrow$ 39 & 否决（一项更差） \\
铺装门（去掉障碍项） & 是（可用率 0.07--0.52 $\rightarrow$ 0.36--0.90） & 6.85 $\rightarrow$ 14.1 m & 0.966 $\rightarrow$ 0.950 & 23 $\rightarrow$ 66 & 否决 \\
交付感知臂 & 否（无分离） & --- & 0.979 $\rightarrow$ 0.995 & 36 $\rightarrow$ 49 & 否决 \\
铺装重定中（原提议） & 未实施 & --- & --- & --- & 被既有测量否定 \\
\bottomrule
\end{tabular}
\end{table*}
\textbf{车道为什么被丢弃。} 铺装门在车道中心不落在观测到的可行驶表面上时撤销感知车道。对门的证据做插桩（两次诊断运行，415 个已结算帧）显示：在 202 个被撤销帧中，\textbf{全部 202 个}都是被障碍项撤销的（约 12 个观测样本中中位数 6 个障碍单元），而「已观测但在可行驶掩码之外」的样本数在其中 192 帧里为零。障碍层由神经鸟瞰图头填充，因此在这段路上，「在铺装上」门实际上成了对车道中心的障碍预测否决。补充材料给出证据与可用率漏斗：在基线配置下，415 个已结算帧中只有 116 帧（28\%）拥有感知车道，只有 79 帧（19\%）拥有规划路径。
\subsection{实时行为与一处指标缺陷}
\textbf{控制节拍。} 逐帧遥测显示感知与规划节拍耗时 377--432 ms，其中 355 ms 是传感器读取（相机环 139--162 ms，距离图像 216--239 ms），11--13 ms 是规划；神经头耗时 3--5 ms。渲染器本身约 50 fps（每帧 20 ms），因此仿真不是瓶颈：控制环按设计以墙钟 2 Hz 步进，且无论如何都受节拍限制在约 2.5 Hz。节拍之间有一个子步路径重跑控制器，每个节拍最多 15 次子步迭代，但它的前提是已经选出路径：在停滞状态下（无路径）不运行子步，指令约每 0.5 s 发出一次；而在路径存在的臂中，每帧 12 个子步带来约每秒 31 次指令发出。这些都是墙钟上的指令发出率，不是关于控制带宽的声明，也不是关于每条指令所依据的传感器数据年龄的声明。表 IV 汇总。实际后果是：车辆被感知到的「卡顿」是路径可用率阻塞导致的控制节拍产物，而不是渲染或算力故障。
\begin{table*}[!t]
\centering
\caption{各配置的控制节拍：节拍耗时、子步数与指令发出率（墙钟）。}
\begin{tabular}{@{}llllll@{}}
\toprule
运行 & 节拍（ms） & 传感器读取（ms） & 渲染帧（ms） & 每帧子步 & 有效指令率 \\
\midrule
基线验收（无路径） & 432 & 401 & 20 & 0 & 约 2.2 Hz \\
车道门因子臂 & 377 & 355 & 20 & 1（最大 5） & 约 5 Hz \\
路径重定中臂 & 403 & 约 370 & 20 & 12 & 约 31 Hz \\
\bottomrule
\end{tabular}
\end{table*}
\textbf{安全读数规则中的一处缺陷。} 用于裁决每个驾驶因子的安全项是车体中心穿越的帧\emph{计数}，而因子臂的行进距离可能是对照臂的两倍。补充材料量化了这一后果：同一臂在每米穿越率减半的同时，可以显示出更大的原始计数。这就是为什么在前两个裁决之后，我们在原始最大值之外另行预注册了按距离的读法（每 100 m 穿越数）；此前的裁决按原样保留，我们把该规则缺陷作为一项方法学结果标出，而不是回头修改它。
\section{讨论}
\textbf{哪些是可推广的。} 本计划中有五件产物似乎可以在具体技术栈之外复用。(i) \textbf{带强制执行凭证的来源阶梯}：为每个标签指明真值级别，并拒绝让命令行字符串把它提升一级，使「这个数字对照的是哪种真值」成为一个可回答的问题；(ii) \textbf{计数契约}：为每个比率报告其分母，并把帧级集合与实例级集合分开，消除了一整类「数字好看、池子为空」的错误，并使逐场景的失败可见，而不是被平均掉。(ii) \textbf{把协议版本化，并把哈希绑定到封存集上}：该机制成本很低，却捕获了一处真实的错标缺陷；没有它，在一套定义下产生的评价就会被归到另一套定义名下。(iii) \textbf{对生成标签做外观认证}：要求被声明为线的像素可从图像中找回，把「标注噪声」从一个调参问题变成一项准入决策。(iv) \textbf{预注册的单因子 A/B 与固定读数规则}，包括明确拒绝把未测量量当作通过；正是这一点让五个负结果得以作为结果而不是失败被报告出来。
\textbf{哪些没有奏效。} 范围缺口无法只靠工程手段弥合，而原因值得说准：约 30\% 的标注线像素不满足采纳定义所用的外观谓词，而该谓词只读 RGB 亮度与色度、不计算与周围铺装的空间对比度，因此这一测量所确立的是「这些像素不满足\emph{这个}谓词」。它们对人是可见的、能否从上下文或另一种模态中找回、以及是否可学，都是本工作没有回答的开放问题------强形式的说法（「不可见，因而不可学」）并不被测量支持。被支持的是算术：漆范围保留开发池 69\% 的标注线像素，而保留下来的像素被召回的比率高于被移除的像素（同一批检查点上 0.813 对 0.736），因此把召回条件化在这个谓词上并不是白得的改进，必须始终与完整标签读法并列报告。在被移除的像素上训练，目前是未测试，而不是已被否定。
离线门也不会自动迁移到闭环验收：同一套通过全部离线门的感知，仍让车辆无法移动，其原因（静止的车体扫掠死锁）没有任何离线指标能表达。最后，可用率与安全沿定义轴相互权衡，这种权衡无法被单个标量概括：在我们已有的四次运行中，采纳定义伴随「未观察到压中心线」与「车道来源可用率相对下降 66\%」同时出现。
\textbf{对评测实践的启示。} 有三种做法本会显著改变我们的结果。报告指标的同时报告其\emph{定义}，并把定义做成版本化产物。对安全计数报告按距离或按机会归一化的结果，因为运动量会污染原始计数。以及，运行一次组成与开发池不同的独立一次性确认；我们那次确认的负结果是整个计划中信息量最大的单项测量。
\begin{table}[!t]
\centering
\caption{被评测各臂的配置，读自交付候选的 checkpoint 元数据。所有臂共用这套配置；组成臂只在训练池组成上不同，单因子臂只在所命名的那个开关上不同。}
\begin{tabular}{@{}ll@{}}
\toprule
项 & 取值 \\
\midrule
结构 & `SegUNet`，轻量 U-Net \cite{ronneberger2015unet}：三个编码块（3x3 卷积、批归一化、2x 最大池化）、带跳跃连接的转置卷积解码器、1x1 头；通道宽度系数 1.0 \\
参数量 & 834,931（由交付 checkpoint 统计） \\
输入 & 536 x 403 RGB 帧 \\
类别 & background、asphalt、line \\
类别权重 & 0.147 / 1.0 / 40.0（线类别加权） \\
损失 & 加权交叉熵 + 线通道的 soft-clDice（权重 1.0）+ Tversky 项 \cite{salehi2017tversky}（alpha 0.3、beta 0.7、权重 1.0）；第 V-C 节被否决的那一臂只改 Tversky 权重 \\
优化器 & Adam（beta 0.9/0.999，无权重衰减），lr 1e-3 \\
学习率计划 & CosineAnnealingLR，T_max 480 步，按迭代步进 \\
批量 & 4 \\
训练预算 & 480 步（每轮 40 步；含续训共记录 12 轮） \\
采样 & 配额采样器覆盖全池；每轮验证比例 0.2（交付臂 158 训练帧 / 29 验证帧） \\
精度 & 启用 AMP；请求确定性算法（CUDA 交叉熵没有确定性实现，这一点如实记录而非隐藏） \\
初始化 & 所有臂从同一个 checkpoint 出发（`t14_e1_promotable_20260927/baseline/seed42`，sha16 `b3a4dd5a82dd4096`）；该 checkpoint 自身的初始化谱系\textbf{未完整记录}（`provenance_complete = False`），因此它之上的谱系无法由存储的元数据重建 \\
硬件 & NVIDIA RTX 5070；torch 2.12.0.dev+cu128、CUDA 12.8、Python 3.10.11 \\
匹配（冻结） & 候选与引擎真值的横向距离在 0.8 m 容差内即算匹配；参考可判要求候选\textbf{自身所在侧}有标注像素；角色一致率比较候选的左右角色与引擎线的角色 \\
汇总方式 & 逐帧累加整数计数、比率只算一次（绝不平均各目录的比率）；分母为空记 UNKNOWN 并附原因，不写 0 \\
数据池 & 人工修订基座 5 目录 / 75 帧；开发池 9 个去重场景 / 117 帧；有限类别池 24 帧；封存的道路不相交留出集 200 帧 \\
\bottomrule
\end{tabular}
\end{table}
\section{局限与效度威胁}
\textbf{不主张实例级召回。} 计数契约报告的覆盖率、身份率与角色都是「进入评价的候选」上的条件比率；感知没有产生候选的帧上发生的漏检只进入像素级召回。需要实例召回的读者应按契约定义（第 III-C 节）理解，而不要从 $|R|/|C|$ 反推一个实例召回。
\textbf{人工基座很小。} 人工修订基座规模偏小（75 个训练帧），人工修订评价包覆盖 9 个包（117 帧），因此场景多样性有限；封存的一次性集合是引擎认证而非人工修订的，因此其数字与人工修订的验收数字不能直接比较。一个更大的人工修订留出集会强化定义性论证，我们将其留给未来工作。在所报告的轮次中，本计划刻意\textbf{不新增人工标注}，这同时限定了成本与覆盖面。
\textbf{单一仿真器、单一地图，而且原因是被量出来的。} 所有感知数据来自同一个仿真器，感知结果来自同一张地图，因为引擎导出的线真值只在该地图上稠密。我们没有假设、而是实测了替代方案：本安装的二十张地图里只有三张声明了类标线贴花材质；在其中 `west_coast_usa` 上生成的受控批次（六个场景、用该地图自己的材质、投影自动标定）\textbf{6/6 全部隔离}，因为引擎的标注通道对生成线报告\textbf{零}线类像素；同一张地图上的原生采集（30 帧、正向相机）线类像素合计 803（约每帧 27），而 `utah` 的正向视角（77 帧）为 0------比管线认证门所需的密度低两个数量级。因此第二张地图需要的是\textbf{另一条真值通道}（在该地图上做标注，或改引擎侧的类别映射），而不是重跑生成器；生成器本身已具备条件（现按地图取材质档案）。定义性结果应按「在一张地图上测得」来读；其方向（而非幅度）已在第二个模型上复现（第 V-B 节）。驾驶结果同样来自一段城镇道路。
\textbf{六个种子，且各臂种子数不等。} 验收使用六个种子；组成臂按各自种子数报告（基础 1、4$\times$ 6、6$\times$ 1、8.5$\times$ 6），而 6$\times$ 剂量点仅依赖被交付的那一个种子。因此剂量结果不是效应量：等种子对照已经做过（表 VI），三档可重评的剂量不可区分，而已记录的 8.5$\times$ 值高于 6$\times$；8.5$\times$ 的检查点已删除，该臂无法重评。同一池内的帧计数高度相关，因此我们报告的变化单位是种子，不把像素当作独立试验。
\textbf{留出集已被消费。} 一次性确认按设计只执行了一次，该集合不得再被读取；未来任何主张都需要重新封存一个集合。实例级探针未在其上运行，仍为 UNKNOWN。
\textbf{产物保留。} 逐轮次训练快照与被取代的运行目录已在一次存储清理中删除；每份判定文件、清单、逐帧遥测、封存帧与交付检查点都被保留，因此本文全部结果与图都可由保留的产物复现，但部分中间训练状态不能。
\textbf{驾驶结果来自研究型技术栈。} 闭环结果描述的是我们的技术栈，而不是一般意义上的自动驾驶；验收门是项目内部目标，其阈值在运行前固定且从未放宽。
\section{结论}
我们试图让自动驾驶的车道感知能够立足在一个规模小但凭证明确的真值基础上被度量，并让它的评价在定义上保持诚实。一条带凭证的来源阶梯，为其引擎一级设置四道认证门，使机器生成的增量可信，而人工基座始终保持小而未被改动；一份计数契约为每个报告的比率给出显式分母；一套版本化、哈希绑定的协议让「线」的定义成为产物而非假设。由此得到的测量恰好在令人不适之处最有信息量：「线」的两套自然定义在同一批数据上各自失败在不同的门上，而这一分歧相当于三分之一的标注像素；采纳定义通过全部离线门，并在闭环中消除全部中心线穿越，代价是可用率减半；一次性独立确认没有复现开发集数字，我们把它作为该检验的头条结果报告。闭环验收门仍未通过，而我们把阻塞点定位到一个可测量的死锁机制，而不是调参不足。每一个因子（包括五个负因子）都经过预注册，并按事前固定的规则裁决；我们认为这种纪律，以及它所产生的那些已记录的负结果，是本文的主要方法学贡献。
\section*{可复现性与数据可得性}
全部结果都由已记录产物重新生成：逐运行记分卡、逐帧遥测、运行清单、门判定文件与封存集记录。两个脚本由这些文件重新生成本文与补充材料的每一张图；协议哈希、封存集摘要与一次性确认记录都印在图中，读者可据此核验数字出自哪套定义与哪批数据。训练池、检查点、封存集与已删除产物的审计都保留在项目仓库中；存储清理期间删除的逐轮次快照列在一份机器可读的审计文件中。
\section*{术语表}
\begin{table}[!t]
\centering
\begin{tabular}{@{}ll@{}}
\toprule
符号 & 含义 \\
\midrule
$P$ & 至少携带一个车道候选的帧 \\
$C$ & 抽取出的候选实例 \\
$C_{out}$ & 落在 $P$ 之外帧上的候选 \\
$R$ & 标注中可用的参考实例 \\
$M$ & 匹配实例，$M \subseteq R$ \\
$L$ & 横向位置一致的实例，$L \subseteq M$ \\
$A$ & 左右角色一致的实例，$A \subseteq L$ \\
覆盖率 & $ & R & / & C & $（实例级） \\
身份率 & $ & M & / & R & $（实例级） \\
角色 & $ & A & / & L & $（实例级） \\
精度、召回 & 像素级，分标签范围与漆范围报告 \\
\bottomrule
\end{tabular}
\end{table}
\begin{table}[!t]
\centering
\begin{tabular}{@{}ll@{}}
\toprule
缩写 & 含义 \\
\midrule
FSD & 全自动驾驶（感知主导的驾驶栈） \\
BEV & 鸟瞰图 \\
IoU & 交并比 \\
TTC & 碰撞时间 \\
SWA & 随机权重平均 \\
PLC & 漆线居中校正器 \\
A/B & 交替的对照／因子比较 \\
UNKNOWN & 该量未被测量；永不读作通过 \\
\bottomrule
\end{tabular}
\end{table}
\section*{未解决的问题与未完成的工作}
本计划刻意把未完成部分与结果以同等详细程度报告；本节逐条列出，并附上定位每一项的证据与具体的下一步实验。
\textbf{1. 闭环验收门尚未通过。} 验收配置的四次运行在基准硬目标上为 0/4（停滞比例 0.93--1.00，120 s 内行进 0.9--14.5 m，距目标剩余 76.5--88.8 m），而已测量的安全项保持（0 碰撞、0 路外帧、0 倒车）。阻塞点被定位为一处静止死锁：一辆偏离中心 0.86 m 的静止车辆，其膨胀车体扫掠在 2.50 m（中位）内穿越车道边界，因此规划器生成的路径被拒绝，车辆无法回到中心（第 V-E 节，图 7 与补充材料）。记录显示的是「已生成的路径被拒绝」，不是「不存在可行路径」；要把两者分开，需要在冻结位姿上做受控扫描，并存下感知边界、候选路径、扫掠几何与最终指令（可行与不可行两类正反例）。\emph{下一步实验}：三个曾经生效的因子（有界保持窗口、路径重定中、车道门）都是单独裁决的；自然的继任者是一次预注册的\textbf{组合}臂对最佳单因子臂的比较，外加计划层才能决定的事项------规划穿越的硬阈值与铺装门的语义，二者都是安全规则。
\begin{figure*}[!t]
\centering
\includegraphics[width=\textwidth]{figures_zh/fig5_deadlock_anatomy.png}
\caption{死锁解剖：规划穿越恰好落在 4 m 阈值之内 (a)，且在这些帧中车辆 96\% 时间静止 (b)。}
\label{fig:deadlock}
\end{figure*}
\textbf{2. 横向参考质量是感知侧的下一根杠杆。} 规划路径距其自身车道参考的中位偏差为 0.773 m，而两层的参考完全一致（差距 0.0 m），有界居中校正器只在存在线证据处生效（805 个已结算帧中的 39 帧）。\emph{下一步实验}：提高线证据可用率（更强的线头，或不依赖稀疏线类别的融合），或把居中前移，使路径直接在车道参考上规划，而不是事后校正。
\textbf{3. 剂量响应：四档中已做三档，且未能复现。} 配对种子对照已经完成（第 V-C 节、表 VI）：六种子下 0$\times$／4$\times$／6$\times$ 不可区分。剩下的是 8.5$\times$ 臂（其检查点在存储清理中删除）与 8.5$\times$ 以上的剂量，两者都需要在冻结协议下\textbf{重新训练}而不是重评。若要主张剂量效应，需要每档等种子、并在不止一张地图上做，且以种子作为变化单位。
\textbf{4. 控制节拍。} 感知与规划节拍耗时 377--432 ms，其中 355 ms 是传感器读取（相机环加距离图像），11--13 ms 是规划；控制环按设计以 2 Hz 步进，并受节拍限制在约 2.5 Hz。子步以最高 15 Hz 平滑间隔，但前提是已选出路径，因此停滞状态恰好是子步不运行的状态。\emph{下一步实验}：更廉价的传感器读取（相机环分辨率或轮转挂载、在既有保活机制下复用距离图像），作为单因子测量，读数同时取节拍成本与驾驶硬目标。
\textbf{5. 安全指标的规格。} 驾驶因子是用帧计数最大值裁决的，而因子臂的行进距离可能是对照臂的两倍；补充材料量化了同一臂如何在每米穿越率减半的同时显示更大的原始计数。新轮次已另行预注册按距离的读法；此前的裁决按记录保留，该规则缺陷作为方法学结果报告，而不是被抹掉。
\textbf{6. 定义性缺口无法靠工程手段弥合。} 约 30\% 的标注线像素是无局部对比度的暗色漆；没有任何依赖局部对比度的模型能找回它们，而在其上训练等于教模型产生幻觉（已记录的负结果）。\emph{可选路径}：显式的约定变更（把它们从召回分母中排除，正如漆范围所做的）、增加一种模态，或引入道路几何先验------每一条都改变了定义，因此都必须作为定义变更来预注册。
\textbf{7. 真值来源与留出容量。} 人工修订基座很小（75 个训练帧；9 个评价包、117 帧），而封存的一次性集合是引擎认证的且已被\textbf{消费}：任何进一步确认都需要重新封存。\emph{下一步实验}：更大的人工修订基座，以及一个重新封存的人工修订留出集------这也会让确认的数字得以与验收数字在同一真值级别上比较。
\textbf{8. 第二张地图：已尝试，被真值通道阻断。} 定义性发现是在一张地图上测得的。我们不只是提议、而是尝试了第二张地图，阻塞点现在是量出来的：引擎导出的线真值在本安装的二十张地图里只在一张上存在。在 `west_coast_usa` 上，我们用该地图自己的材质与自动标定生成了受控批次，结果 6/6 全部隔离（标注通道对生成线报告零线类像素）；同一张地图的原生采集每帧约 27 个线类像素，`utah` 则为零。剩下两条路，都是真正的工作而不是重跑：为新地图做标注（这会打破本计划一直保持的「不新增人工标注」），或扩展引擎侧的类别映射，让第二张地图的标线被标注为线。在其中一条完成之前，范围敏感性应按「在一张地图上测得」来对待，其方向已在第二个模型上复现。
\textbf{9. 更复杂几何下的实例词表。} 计数契约的角色词表（左／右、近／远）是为单一车行道上的车道标线定义的；交叉路口、3D 车道估计与多车道角色约定都尚未定义，而基线定义下身份门的失败模式恰恰是角色／实例歧义，而不是像素误差。
\section*{声明}
\textbf{基金。} 本研究未获得外部资助。
\textbf{利益冲突。} 作者声明不存在利益冲突。
\textbf{作者贡献。} 作者为唯一作者，承担全部角色：概念构思、方法学、软件、验证、形式分析、调查、数据管理、写作（初稿）、写作（评审与编辑）、可视化、监督。
\textbf{伦理。} 不适用：本工作使用商用驾驶仿真器，不涉及人类受试者。本文提及的人工标注由作者本人对自己录制的仿真数据完成。
\textbf{AI 工具的使用。} 在实现与文档撰写过程中使用了 AI 编程助手（DeepSeek 与 OpenCode）作为开发工具；全部实验设计、测量、裁决与结论均由作者产出并核验，本文的每一项定量主张都由下文所列脚本从已记录产物重新生成。
\textbf{数据与代码可得性。} 公开的是\textbf{代码与稿件源}：`https://github.com/Qiongkura/BeamNG-autopilot`（分支 `fix/round3-hardening-20260921`）包含评测栈、协议与计数契约定义、图生成脚本、稿件 Markdown、`build_paper.py` 与 `references.bib`。\textbf{不公开的是记录产物}：逐运行记分卡、逐帧遥测、运行清单、门判定文件、封存集记录与检查点都位于作者工作副本的 `logs/` 下，该目录未被仓库跟踪（忽略规则排除）。本文每一个数字都由附录 A 所列脚本从这些本地产物重新生成，`CLAIMS_20261007.md` 逐条列出每项主张所依据的输入文件、协议、种子集合与范围，因此这套对应关系可以逐文件核对。产物的公开归档\textbf{尚不可用}，需要者请联系作者。存储清理期间删除的逐轮次训练快照与被取代的运行目录列在机器可读审计文件（`logs/_cleanup_20261006.txt`）中，这也是第 V-C 节等种子剂量对照只覆盖四档中的三档的原因。
\section*{致谢}
作者感谢 BeamNG.tech 团队提供仿真器及其脚本接口，并衷心感谢父母与朋友的支持。作者同时感谢 DeepSeek 与 OpenCode 的开发者，其编程助手在本工作中作为开发工具使用。
\section*{附录 A：各结果节对应的已记录产物}
\begin{table}[!t]
\centering
\begin{tabular}{@{}ll@{}}
\toprule
结果节 & 产物 \\
\midrule
V-A 验收 & `r2_verdict_v7_20261005.json`, `r2_verdict_v8_ADOPTED_20261005.json` \\
V-B 定义性缺口 & `r2_verdict_v8_ADOPTED_20261005.json`（逐种子标签与漆范围召回） \\
V-B 边界图 & `r2_verdict_v8_paint_{nodilate_mean,dil1b,dil2}_20261005.json` \\
V-B 外观门 & `dev_pixel_{base,appgate}_20261004.json`, `r3b_pixel_{lat3,appgate}_20261004.json` \\
V-B 扫描 & `dev_scan_{lateral,parallel}_20261004.json`, `r3b_scan_{lateral,parallel}_20261004.json` \\
V-B 身份范围 & `t16_dual_scope_10seed.json` \\
V-C 剂量--效应 & `dev_full_arms_20261005.json`, `dev_full_dose_arms_s4{3..7}.json` \\
V-C 被否决因子 & `r2_verdict_swa_20261005.json`, `dev9_beta08_v8d1_20261005.json` \\
V-D 确认 & `final_set_v5_20261005/seal/{final_set_seal.json,confirmation_base6x-seed42.json,final_set_ledger.jsonl}` \\
V-D 组成审计 & `final_pool_audit_20261005.json`, `collect_isolation_20260927.json` \\
V-E 驾驶 & `logs/fsd_benchmark/{scorecard,manifest,town}_*.json`（每次 A/B 的全部运行） \\
V-F 计时 & `timing_retest_20260927.json`，逐帧遥测的 `frame_ms`/`tick_ms` 字段 \\
IV 池筛选 & `t16_r3_density_dev.json`, `t16_r3_pool{,b}_density_scenes.json`, `e1_negative_pool_20260927.json` \\
IV 成本台账 & `gpu_ledger_machine.json` \\
\bottomrule
\end{tabular}
\end{table}
\section*{参考文献}
（由 `references.bib` 自动生成。）

\begin{thebibliography}{99}
\bibitem{neven2018lanenet} Neven, Davy, De Brabandere, Bert, Georgoulis, Stamatios, Proesmans, Marc, Van Gool, Luc, “Towards End-to-End Lane Detection: An Instance Segmentation Approach,” in IEEE Intelligent Vehicles Symposium (IV), 2018, doi: 10.1109/IVS.2018.8500547.
\bibitem{pan2018scnn} Pan, Xingang, Shi, Jianping, Luo, Ping, Wang, Xiaogang, Tang, Xiaoou, “Spatial As Deep: Spatial CNN for Traffic Scene Understanding,” in AAAI Conference on Artificial Intelligence, 2018, doi: 10.1609/aaai.v32i1.12301.
\bibitem{zheng2022clrnet} Zheng, Tu, Huang, Yifei, Liu, Yang, Tang, Wenjian, Yang, Zheng, Cai, Deng, He, Xiaofei, “CLRNet: Cross Layer Refinement Network for Lane Detection,” in IEEE/CVF Conference on Computer Vision and Pattern Recognition (CVPR), 2022, doi: 10.1109/CVPR52688.2022.00097.
\bibitem{chen2026clrnetlight} 陈淑慧, “CLRNet车道线检测算法的一种轻量化方法,” 软件工程, 29, no. 9, 2026. [In Chinese]
\bibitem{huang2026bspline} 黄炜, 等, “基于B样条曲线的端到端车道线检测方法,” 长安大学学报(自然科学版), 46, no. 3, 118--129, 2026, doi: 10.19721/j.cnki.1671-8879.2026.03.008. [In Chinese]
\bibitem{jia2026gnn} 贾子厚, 等, “基于GNN-Transformer模型的车道线检测方法,” 工程设计学报, 33, no. 3, 2026. [In Chinese]
\bibitem{li2026xca} 李睿, 敖银辉, 陈新盛, 李桂伸, “基于Transformer和XCA注意力的车道线检测方法,” 电子测量技术, 49, no. 7, 2026, doi: 10.19651/j.cnki.emt.2519485. [In Chinese]
\bibitem{zhu2024parallel} 朱威, 欧全林, 洪力栋, 何德峰, “基于图像序列的车道线并行检测网络,” 模式识别与人工智能, 34, no. 5, 434--445, 2021. [In Chinese]
\bibitem{hu2026multiscale} 胡鹏, 刘平, 胡遥, “基于多尺度特征聚合的车道线检测方法,” 汽车工程学报, 2026. [In Chinese]
\bibitem{shuai2026deeplab} 帅强, 等, “基于改进DeeplabV3+模型的车道线检测方法,” 太原科技大学学报, 47, no. 3, 208--214, 2026, doi: 10.3969/j.issn.1673-2057.2026.03.006. [In Chinese]
\bibitem{chen2018deeplab} Chen, Liang-Chieh, Zhu, Yukun, Papandreou, George, Schroff, Florian, Adam, Hartwig, “Encoder-Decoder with Atrous Separable Convolution for Semantic Image Segmentation,” in European Conference on Computer Vision (ECCV), 2018, doi: 10.1007/978-3-030-01234-2_49.
\bibitem{lu2026symmetry} 陆军, 张澍森, 王磊, “复杂场景下基于对称先验与关键点预测的车道线检测,” 智能系统学报, 2026, doi: 10.11992/tis.202602011. [In Chinese]
\bibitem{chen2025threed} 陈小海, 等, “深度卷积神经网络架构下智能车3D车道线检测研究,” 现代电子技术, 48, no. 24, 2025. [In Chinese]
\bibitem{tusimple2017} TuSimple, “TuSimple Lane Detection Challenge,” 2017. [Dataset]
\bibitem{cordts2016cityscapes} Cordts, Marius, Omran, Mohamed, Ramos, Sebastian, Rehfeld, Timo, Enzweiler, Markus, Benenson, Rodrigo, Franke, Uwe, Roth, Stefan, Schiele, Bernt, “The Cityscapes Dataset for Semantic Urban Scene Understanding,” in IEEE Conference on Computer Vision and Pattern Recognition (CVPR), 2016, doi: 10.1109/CVPR.2016.350.
\bibitem{borkar2012groundtruth} Borkar, Amol, Hayes, Monson, Smith, Mark T., “A Novel Lane Detection System With Efficient Ground Truth Generation,” IEEE Transactions on Intelligent Transportation Systems, 13, no. 1, 365--374, 2012, doi: 10.1109/TITS.2011.2173196.
\bibitem{mccall2006video} McCall, Joel C., Trivedi, Mohan M., “Video-Based Lane Estimation and Tracking for Driver Assistance: Survey, System, and Evaluation,” IEEE Transactions on Intelligent Transportation Systems, 7, no. 1, 20--37, 2006, doi: 10.1109/TITS.2006.869595.
\bibitem{izmailov2018swa} Izmailov, Pavel, Podoprikhin, Dmitrii, Garipov, Timur, Vetrov, Dmitry, Wilson, Andrew Gordon, “Averaging Weights Leads to Wider Optima and Better Generalization,” in Conference on Uncertainty in Artificial Intelligence (UAI), 2018. [arXiv:1803.05407]
\bibitem{ronneberger2015unet} Ronneberger, Olaf, Fischer, Philipp, Brox, Thomas, “U-Net: Convolutional Networks for Biomedical Image Segmentation,” in Medical Image Computing and Computer-Assisted Intervention (MICCAI), 2015, doi: 10.1007/978-3-319-24574-4_28.
\bibitem{salehi2017tversky} Salehi, Seyed Sadegh Mohseni, Erdogmus, Deniz, Gholipour, Ali, “Tversky Loss Function for Image Segmentation Using 3D Fully Convolutional Deep Networks,” in Machine Learning in Medical Imaging (MLMI), MICCAI Workshop, 2017, doi: 10.1007/978-3-319-67389-9_44.
\end{thebibliography}
\end{document}
